% !PRINTED 178
\weberpage{178}
A large class of subgroups of the group $\mathfrak{L}$ with finite
index is formed by the so-called \gesperrt{congruence groups}, which, if
$m$ is any integer, are characterised by the four congruences
\[
\begin{pmatrix} \alpha, & \beta \\ \gamma, & \delta \end{pmatrix} \equiv \begin{pmatrix} 1, & 0 \\ 0, & 1 \end{pmatrix} (\mathrm{mod}\ m)
\]
The modular functions belonging to these groups are
 called \gesperrt{congruence moduli} of the $m$th level. For $m = 2$
this is the group $\mathfrak{A}$ (\S\ 36). We shall make use of the
following theorems:

1. \gesperrt{By \S\ 52 the modulus $\varkappa^2$ is a modular function belonging
 to this group}; and since by \S\ 36 all these transformations
can be composed from the two
\[
\begin{pmatrix} 1, & 0 \\ 2, & 1 \end{pmatrix}, \qquad \begin{pmatrix} 1, & 2 \\ 0, & 1 \end{pmatrix}
\]
we can state the theorem:

\gesperrt{Every modular function which remains unchanged under the two substitutions
}
\[
(\omega,\ \omega+2), \qquad \left(\omega,\ \frac{\omega}{1+2\omega}\right)
\]
\gesperrt{is a rational function of $\varkappa^2$.}

2. The function $\varkappa^2\varkappa'^2$ belongs to the group $\mathfrak{A}'$ formed
jointly from the first and second classes of \S\ 36.
This group can be derived by repetition of the two transformations
\[
\begin{pmatrix} 1, & 0 \\ 2, & 1 \end{pmatrix}, \qquad \begin{pmatrix} 0, & 1 \\ -1, & 0 \end{pmatrix}
\]
and we can therefore state the theorem:

\gesperrt{Every modular function which remains unchanged under the two substitutions
}
\[
(\omega,\ \omega+2), \qquad \left(\omega,\ -\frac{1}{\omega}\right)
\]
\gesperrt{is a rational function of $\varkappa^2\varkappa'^2$.}
From this there follows further:

3. \gesperrt{Every modular function which remains unchanged under $(\omega,\omega+2)$
and changes its sign under $\left(\omega, -\dfrac{1}{\omega}\right)$ is the
product of $(\varkappa'^2 - \varkappa^2)$ with a rational function
of $\varkappa^2\varkappa'^2$.}
